% SGA 3, Expose VI_A, Loop 1 English reconstruction.
% Scope: section 3.3 in full, through Corollary 3.3.3;
% section 4 is excluded.
% Authority: Polo--Gille Exp6A-13oct24.pdf, component-local pages 25--26
% (combined-reader pages 328--329; printed pages 317--318).
% OCR and the Reinhold English Markdown are comparison witnesses only.
% The exact sequence is an exact 180-dpi crop of the controlling PDF:
% local p. 25, x=500,y=1546,w=540,h=78 on a 1488x2105 render.

\medskip
\phantomsection\label{par:sga3-VIA-3.3}%
\noindent\textbf{3.3.---}%
\footnote{\label{note:sga3-VIA-42}编者注: 此小节为新增内容.}
在\hyperref[thm:sga3-VIA-3.2]{定理~3.2} 中,
当 \(u:F\to G\) 为单态射时, 可省去 \(G\) 平坦的假设.
第 VI\textsubscript{B} 篇, 注~9.3(b),
以及 \textup{[Ray67a]}, 例~a)\SGARefLink{thm:sga3-VIA-3.2:item:i}{(i)}, 第~82~页,
都提到了此推广. \textup{[An73]}, \SGARefLink{sec:sga3-VIA-4}{定理~4} 给出的证明,
由\hyperref[thm:sga3-VIA-3.2]{定理~3.2} 和下述 Grothendieck 定理得到
(该定理见于 \textup{[Ray67a]}, 定理~1(ii),
证明见 \textup{[DG70]}, \S~III.2, \SGARefLink{prop:sga3-VIB-7.1}{7.1}).
若 \(X\) 为概形, \(R\) 为 \(X\) 上的等价关系,
用 \(\widetilde{X/R}\) 表示 \(X\) 对 \(R\) 的 \((\mathrm{fppf})\) 商层
(参见\SGARefLink{sga3:iv:proposition:4:4:9}{第 IV 篇, 4.4.9}).

\medskip
\phantomsection\label{thm:sga3-VIA-3.3.1}%
\noindent\textbf{定理 3.3.1 (Grothendieck).}
\textit{设 \(A\) 为环, \(X\) 为 \(A\)-概形,
\[
  d_0,d_1:R\rightrightarrows X
\]
为 \(X\) 上的 \(A\)-等价关系, 其中每个 \(d_i\) 都忠实平坦且有限呈示.
设 \(X_0\) 为 \(X\) 的由幂零理想定义的饱和闭子概形,
\(R_0\) 为 \(R\) 在 \(X_0\) 上诱导的等价关系.
若 \((\mathrm{fppf})\) 商层 \(\widetilde{X_0/R_0}\) 可由 \(A\)-概形表出,
则 \(\widetilde{X/R}\) 也可表.}

证明见 \textup{[DG70]}, \S~III.2, \SGARefLink{prop:sga3-VIB-7.1}{7.1}.
现在回到 \(A\) 为 Artin 局部环的情形.
设 \(u:F\to G\) 为局部有限型 \(A\)-群之间的拟紧态射,
再假设 \(u\) 为单态射.
则由\SGARefLink{prop:sga3-VIA-2.5.2}{2.5.2}, \(u\) 为闭浸入.

现在可以陈述\hyperref[thm:sga3-VIA-3.2]{定理~3.2} 的如下变式.

\medskip
\phantomsection\label{thm:sga3-VIA-3.3.2}%
\noindent\textbf{定理 3.3.2.}
\textit{设 \(A\) 为 Artin 局部环, \(G\) 为局部有限型 \(A\)-群,
\(F\) 为 \(G\) 的在 \(A\) 上平坦的闭子群.}%
\footnote{\label{note:sga3-VIA-43}编者注: 关于 \(F\) 不平坦且 \(\widetilde{F\backslash G}\) 不可表的例子,
见 \textup{[DG70]}, \S~III.3, 第~\SGARefLink{prop:sga3-VIB-3.3}{3.3} 号.}
\textit{则:}
\begin{enumerate}
\item[\SGAInlineRefTarget{thm:sga3-VIA-3.3.2:item:i}(i)] \textit{\((\mathrm{fppf})\) 商层 \(\widetilde{F\backslash G}\)
可由分离的局部有限型 \(A\)-概形 \(F\backslash G\) 表出;
此外, 序列}
\begin{center}
  \phantomsection\label{diag:sga3-VIA-3.3.2-exact-sequence}%
  \begin{tikzcd}[column sep=large]
    F\times_A G
      \arrow[r,shift left=1.2ex]
      \arrow[r,shift right=1.2ex]
    &
    G
      \arrow[r,"p"]
    &
    F\backslash G
  \end{tikzcd}
\end{center}
\textit{在所有赋环空间组成的范畴中正合.}

\item[\SGAInlineRefTarget{thm:sga3-VIA-3.3.2:item:ii}(ii)] \textit{态射
\[
  F\times_A G
    \xrightarrow{(\lambda,\operatorname{pr}_2)}
  G\times_{(F\backslash G)}G
\]
为同构, 且 \(p:G\to F\backslash G\) 忠实平坦并局部有限呈示;
故 \(p\) 为关于 \((\mathrm{fppf})\) 拓扑局部平凡的 \(F\)-挠子.}

\item[\SGAInlineRefTarget{thm:sga3-VIA-3.3.2:item:iii}(iii)] \textit{若 \(G\) 在 \(A\) 上平坦 (相应地, 有限型; 相应地, 平滑),
则 \(F\backslash G\) 也如此.}

\item[(iv)]\phantomsection\label{concl:sga3-VIA-3.3.2-iv} \textit{\(X=F\backslash G\) 带有 \(G\) 的右作用,
使对每个 \(g\in G\), \(p(e)\cdot g=p(g)\).
因此 \(X\) 的连通分支在 \(A\) 上有限型、不可约,
且维数均为 \(\dim G-\dim F\).}

\item[(v)]\phantomsection\label{concl:sga3-VIA-3.3.2-v} \textit{此外, 若 \(F\) 为 \(G\) 的正规子群,
则 \(F\backslash G\) 上有唯一的 \(A\)-群结构,
使 \(p:G\to F\backslash G\) 为 \(A\)-群的态射.}
\end{enumerate}

断言\SGARefLink{thm:sga3-VIA-3.3.2:item:i}{(i)}
和\SGARefLink{thm:sga3-VIA-3.3.2:item:ii}{(ii)}
由\SGARefLink{thm:sga3-VIA-3.2}{3.2} 和\SGARefLink{thm:sga3-VIA-3.3.1}{3.3.1} 得到.
由于 \(G\to F\backslash G\) 忠实平坦且局部有限呈示,
断言 (iii) 由 EGA~IV, 2.2.14, 2.7.1 和 17.7.7 得到.
断言 (iv) 和 (v) 将在\hyperref[sec:sga3-VIA-5]{第~5~节}证明.
立即记下如下推论.

\medskip
\phantomsection\label{cor:sga3-VIA-3.3.3}%
\noindent\textbf{推论 3.3.3.}
\textit{设 \(A\) 为 Artin 局部环, \(G\) 为局部有限型 \(A\)-群,
\(H\) 为 \(G\) 的在 \(A\) 上平坦的闭子群.
用 \(p\) 表示态射 \(G\to G/H\), 用 \(\lambda\)
(相应地, \(\operatorname{pr}_1\)) 表示由 \(\lambda(g,h)=gh\)
定义的态射 \(G\times H\to G\) (相应地, 投影 \(G\times H\to G\)).
则对 \(G/H\) 的每个开子集 \(U\),
\[
  \mathcal O(U)
  =
  \bigl\{
    \varphi\in\mathcal O\bigl(p^{-1}(U)\bigr)
    \,\bigm|\,
    \varphi\circ\lambda=\varphi\circ\operatorname{pr}_1
  \bigr\}.
\]
换言之, \(\mathcal O(U)\) 是那些满足如下条件的
\(\varphi\in\mathcal O(p^{-1}(U))\) 所成的集合:
对每个 \(A\)-概形 \(S\) 及每个 \(g\in G(S)\), \(h\in H(S)\),
都有 \(\varphi(gh)=\varphi(g)\).}

事实上, \(p:G\to G/H\) 忠实平坦且局部有限呈示,
从而为 \((\mathrm{fppf})\) 拓扑的覆盖态射;
此结果由\SGARefLink{sga3:iv:remark:3:3:3:2}{第 IV 篇, 3.3.3.2} 得到.
% Independent prose review: source and translation compared paragraph by paragraph through this component; statements, hypotheses, proofs, notes and diagram captions covered; bibliography reserved for its assigned reviewer; source math and reference keys unchanged.
